
%=========== File containing some new commands ============%
%                                                          %
% Copyright (C) ISI - All Rights Reserved                  %
% Proprietary                                              %
% Written by Med Hossam <med.hossam@gmail.com>, April 2016 %
%                                                          %
% @author: HEDHILI Med Houssemeddine                       %
% @linkedin: http://tn.linkedin.com/in/medhossam           %
%==========================================================%

\newenvironment{changemargin}[2]{%
\begin{list}{}{%
\setlength{\leftmargin}{#1}%
\setlength{\rightmargin}{#2}%
}%
\item[]}
{\end{list}}

\makeatletter

%================= front cover variables =================%

\newcommand{\secondAuthor}[1]{\gdef\@secondAuthor{#1}}%
\newcommand{\@secondAuthor}{\@latex@warning@no@line{No \noexpand\secondAuthor given}}

\newcommand{\diplomaName}[1]{\gdef\@diplomaName{#1}}%
\newcommand{\@diplomaName}{\@latex@warning@no@line{No \noexpand\diplomaName given}}

\newcommand{\speciality}[1]{\gdef\@speciality{#1}}%
\newcommand{\@speciality}{\@latex@warning@no@line{No \noexpand\speciality given}}

\newcommand{\proFramerName}[1]{\gdef\@proFramerName{#1}}%
\newcommand{\@proFramerName}{\@latex@warning@no@line{No \noexpand\proFramerName given}}

\newcommand{\proFramerSpeciality}[1]{\gdef\@proFramerSpeciality{#1}}%
\newcommand{\@proFramerSpeciality}{\@latex@warning@no@line{No \noexpand\proFramerSpeciality given}}

\newcommand{\academicFramerName}[1]{\gdef\@academicFramerName{#1}}%
\newcommand{\@academicFramerName}{\@latex@warning@no@line{No \noexpand\academicFramerName given}}

\newcommand{\academicFramerSpeciality}[1]{\gdef\@academicFramerSpeciality{#1}}%
\newcommand{\@academicFramerSpeciality}{\@latex@warning@no@line{No \noexpand\academicFramerSpeciality given}}

\newcommand{\collegeYear}[1]{\gdef\@collegeYear{#1}}%
\newcommand{\@collegeYear}{\@latex@warning@no@line{No \noexpand\collegeYear given}}

\newcommand{\companyName}[1]{\gdef\@companyName{#1}}%
\newcommand{\@companyName}{\@latex@warning@no@line{No \noexpand\companyName given}}

%================== Signatures variables ==================%

\newcommand{\proSignSentence}[1]{\gdef\@proSignSentence{#1}}%
\newcommand{\@proSignSentence}{\@latex@warning@no@line{No \noexpand\proSignSentence given}}

\newcommand{\academicSignSentence}[1]{\gdef\@academicSignSentence{#1}}%
\newcommand{\@academicSignSentence}{\@latex@warning@no@line{No \noexpand\academicSignSentence given}}

%================== Backcover variables ==================%

\newcommand{\arabicAbstract}[1]{\gdef\@arabicAbstract{#1}}%
\newcommand{\@arabicAbstract}{\@latex@warning@no@line{No \noexpand\arabicAbstract given}}

\newcommand{\arabicAbstractKeywords}[1]{\gdef\@arabicAbstractKeywords{#1}}%
\newcommand{\@arabicAbstractKeywords}{\@latex@warning@no@line{No \noexpand\arabicAbstractKeywords given}}

\newcommand{\frenchAbstract}[1]{\gdef\@frenchAbstract{#1}}%
\newcommand{\@frenchAbstract}{\@latex@warning@no@line{No \noexpand\frenchAbstract given}}

\newcommand{\frenchAbstractKeywords}[1]{\gdef\@frenchAbstractKeywords{#1}}%
\newcommand{\@frenchAbstractKeywords}{\@latex@warning@no@line{No \noexpand\frenchAbstractKeywords given}}

\newcommand{\englishAbstract}[1]{\gdef\@englishAbstract{#1}}%
\newcommand{\@englishAbstract}{\@latex@warning@no@line{No \noexpand\englishAbstract given}}

\newcommand{\englishAbstractKeywords}[1]{\gdef\@englishAbstractKeywords{#1}}%
\newcommand{\@englishAbstractKeywords}{\@latex@warning@no@line{No \noexpand\englishAbstractKeywords given}}

\newcommand{\companyEmail}[1]{\gdef\@companyEmail{#1}}%
\newcommand{\@companyEmail}{\@latex@warning@no@line{No \noexpand\companyEmail given}}

\newcommand{\companyTel}[1]{\gdef\@companyTel{#1}}%
\newcommand{\@companyTel}{\@latex@warning@no@line{No \noexpand\companyTel given}}

\newcommand{\companyFax}[1]{\gdef\@companyFax{#1}}%
\newcommand{\@companyFax}{\@latex@warning@no@line{No \noexpand\companyFax given}}

\newcommand{\companyAddressFR}[1]{\gdef\@companyAddressFR{#1}}%
\newcommand{\@companyAddressFR}{\@latex@warning@no@line{No \noexpand\companyAddressFR given}}

\newcommand{\companyAddressAR}[1]{\gdef\@companyAddressAR{#1}}%
\newcommand{\@companyAddressAR}{\@latex@warning@no@line{No \noexpand\companyAddressAR given}}

%============= cmd for inserting blank page =============%
\newcommand\blankpage{%
    \null
    \thispagestyle{empty}%
    \addtocounter{page}{-1}%
    \newpage}

%================ document main language ================%
%\selectlanguage{english}
\selectlanguage{french}

%================== required packages ===================%

\usepackage{tcolorbox}
\usepackage{afterpage}
\usepackage{array,longtable,multirow}% http://ctan.org/pkg/{array,longtable,multirow}
\usepackage{pifont}

\usepackage{pdflscape}
\usepackage{rotating}
%===== Numero et titre courant sur les pages paysage =====%
% Ce sont les deux lignes que la classe ecrit elle-meme (tpl/isipfe.cls:305-306)
% dans son style lscape, et qu'elle a laissees en commentaire. isipfe.cls devant
% rester intact, elles sont reprises ici aux valeurs corrigees. fancyhdr ne permute pas
% l'en-tete et le pied quand lscape permute le texte : le \rotatebox{90} remet
% les glyphes droits une fois la page tournee par /Rotate, et un \raisebox les
% deplace le long du grand cote. Les trois builds, en abscisse LUE :
% build 9, vspace du modele (+12cm / -88mm) = numero a 819, titre a -163 rogne
% hors papier ; build 10, vspace -2,2cm des deux cotes = numero a 755, titre
% inchange ; build 11, raisebox -11,9cm = numero a 481, titre a 270 et entier.
% Le raisebox est lineaire et la case de tete a sa baseline naturelle en
% x = 819,6 : 481 = 819,6 - 338,6 (11,9cm) et 755 = 819,6 - 62,4 (2,2cm). Le
% \vspace, lui, sature des qu'il monte au-dessus de la boite de \headheight
% (+12cm rend 819 comme rien du tout) : d'ou le raisebox, seul reglable au pt.
% Le folio ne peut PAS etre centre sur le grand cote bas, ni meme pose au
% milieu : la legende de la figure occupe toute cette bande. Mesure build 11 :
% la legende est centree sous l'image (centre x = 435), donc sa portee se lit a
% partir de son ancre — la plus longue court de x = 140 a x = 730 (p. 84,
% legende de deux lignes), les plus courtes de 351 a 519, et sa baseline est a
% y = 74, 75 ou 76 quand celle du folio est a y = 68, six points plus bas. Pose
% a x = 481, comme centre a 421, le nombre s'imprime donc dans la legende — 15
% des 25 pages sont signalees par le controle. Il ne reste que le coin :
% x >= 770 pour laisser 40 pt apres la legende la plus longue, et x + 14 <= 797
% pour rester a plus de 1,5 cm du bord de papier (la page lue fait 842 pt de
% large). x = 775 retenu, soit 819,6 - 44,1 pt = -1,55cm. Le folio reste ainsi
% dans le meme bandeau bas que les pages portrait, au coin droit, et la classe
% n'est pas touchee.
% Le titre courant, dans la bande haute (la case de pied permute), est a
% y = 531, au-dessus de l'image (y = 524 au plus) ; son depart x = 270 est le
% meme sur les 25 pages et sa longueur ne depasse pas 200 pt, rien ne sort du
% papier. Sans regle : la classe les annule (cls:303-304) et on la suit.
\fancypagestyle{lscape}{%
  \fancyhf{}%
  \renewcommand{\footrulewidth}{0pt}%
  \renewcommand{\headrulewidth}{0pt}%
  \rhead{\raisebox{-1.55cm}[0pt][0pt]{\rotatebox{90}{\footnotesize \thepage}}}%
  \lfoot{\raisebox{8cm}[0pt][0pt]{\rotatebox{90}{\nouppercase{\leftmark}}}}%
}
\usepackage{wrapfig}
%=========== Ajoutes pour le contenu du rapport (code, URLs, tableaux) ===========%
\newcolumntype{P}[1]{>{\raggedright\arraybackslash}p{#1}}
\usepackage{listings}
\usepackage{url}
\lstdefinestyle{CodePFE}{
  inputencoding=utf8,
  literate={é}{{\'e}}1 {è}{{\`e}}1 {à}{{\`a}}1 {ê}{{\^e}}1 {î}{{\^i}}1
           {ô}{{\^o}}1 {û}{{\^u}}1 {ç}{{\c{c}}}1 {É}{{\'E}}1 {À}{{\`A}}1
           {«}{{\guillemotleft}}1 {»}{{\guillemotright}}1 {—}{{---}}1,
  basicstyle=\ttfamily\footnotesize,
  keywordstyle=\color{isiBlue}\bfseries,
% Une seule listing de 24 lignes dans ce rapport. Declarer « language=Java »
% force listings a parser lstlang1.sty (environ 100 ko de definitions de mots-cles)
% a CHAQUE passe de compilation, soit trois fois pour un seul extrait : c'est le
% poste fixe le plus couteux du document et il fait depasser la limite de 20 s du
% plan gratuit. On garde le gras colore en enumerant ici les seuls mots-cles que
% contient l'extrait ; le rendu de l'extrait est identique.
  morekeywords={public,private,protected,class,interface,extends,implements,import,void,new,return,this,true,false},
  columns=fullflexible,
  breaklines=true,
  breakatwhitespace=true,
  keepspaces=true,
  showstringspaces=false,
  frame=single,
  rulecolor=\color{black!25},
  numbers=left,
  numberstyle=\tiny\color{black!45},
  numbersep=6pt,
  xleftmargin=2em,
  aboveskip=10pt,
  belowskip=10pt
}

%=========== Les identifiants techniques ne debordent plus de la marge ===========%
% « MANAGER\_REJECTED », « AUTO\_REJECTED »… tiennent rarement sur une ligne et
% sont restes bloques, car un souligne interdit normalement la césure. On autorise
% ici un retour à la ligne juste apres chaque souligne, sans trait d'union.
\def\_{\textunderscore\allowbreak}

% Dernier filet : autoriser un leger etirement extra-ligne plutot que de laisser
% une ligne depasser la marge droite.
\setlength{\emergencystretch}{1em}

%=========== Floats : eviter que LaTeX gaspille des pages ===========%
% 16 figures en paysage + 27 en portrait : sans ces seuils, LaTeX refuse les
% pages hautes et repousse les flottants, ce qui cree des trous et rallonge la
% compilation (limite 20 s du plan gratuit Overleaf).
\renewcommand{\topfraction}{0.95}
\renewcommand{\bottomfraction}{0.85}
\renewcommand{\textfraction}{0.05}
\renewcommand{\floatpagefraction}{0.75}
\setcounter{totalnumber}{3}
%=========== Fin floats ======%
